% cover-tikz.tex — Colfax《CuTe 布局的范畴论基础》中文译注版 封面（手绘风，满幅）
% 风格参照 ../cute-zhihu/cover-tikz.tex：Humor Sans + 霞鹜文楷、sketch 手绘边框、
% 层叠卡片、散落 mono 元组纹理、底部色条与落款；主色改 Colfax 紫 + colorx 八色 ramp。
% 版面按 colfax 实际纸张 letterpaper (21.59cm x 27.94cm) 绘制，
% 与 \paperwidth/\paperheight 一致，shipout 铺图时不会拉伸变形。
% 构建：xelatex cover-tikz.tex（standalone，输出单页 PDF cover-tikz.pdf）
\documentclass[tikz,border=0pt]{standalone}
\usetikzlibrary{arrows.meta,calc,decorations.pathmorphing,shadows}
\usepackage{xeCJK}
\usepackage{amsmath}
\setmainfont{Humor Sans}
\setsansfont{Humor Sans}
\setmonofont{Comic Neue}
\setCJKmainfont{LXGW WenKai}
\setCJKsansfont{LXGW WenKai}

\definecolor{paperBg}{RGB}{252,251,247}
\definecolor{cfInk}{RGB}{30,34,40}
\definecolor{cfMuted}{RGB}{120,128,140}
\definecolor{cfViolet}{RGB}{126,99,184}
\definecolor{cfDeep}{RGB}{96,68,148}
\definecolor{warmYellow}{RGB}{232,163,61}
\definecolor{coral}{RGB}{213,84,76}
\definecolor{mint}{RGB}{76,182,149}
\definecolor{colorx1}{RGB}{213,84,76}
\definecolor{colorx2}{RGB}{232,163,61}
\definecolor{colorx3}{RGB}{155,191,90}
\definecolor{colorx4}{RGB}{76,182,149}
\definecolor{colorx5}{RGB}{69,168,200}
\definecolor{colorx6}{RGB}{62,124,192}
\definecolor{colorx7}{RGB}{126,99,184}
\definecolor{colorx8}{RGB}{194,88,158}

\tikzset{
  sketch/.style={decorate, decoration={random steps, segment length=2.5mm, amplitude=0.35mm}},
  card/.style={sketch, rounded corners=6pt, line width=1.2pt, draw=cfInk!70,
    fill=white, fill opacity=0.92, drop shadow={shadow xshift=1.2mm, shadow yshift=-1.6mm, opacity=0.14}},
  obj/.style={draw=cfViolet!65, fill=cfViolet!8, rounded corners=2pt, inner sep=2.5pt,
    text=cfInk, font=\ttfamily\scriptsize}}

% 成书日期取编译当月，勿写死（write-leetcuda-book 硬规则 11）
\newcommand{\builddate}{\the\year 年 \the\month 月}

\begin{document}
\pgfmathsetseed{20260901}  % 固定 random steps 种子，保证重复编译产物一致
\begin{tikzpicture}
% 背景：暖纸白 -> 极淡 Colfax 紫
\shade[top color=paperBg, bottom color=paperBg!92!cfViolet] (0,0) rectangle (21.59,27.94);

% 背景装饰：散落 layout 元组与范畴论符号
\node[rotate=10, text=cfMuted!38, font=\ttfamily\footnotesize] at (2.5,24.4) {(8,16):(1,8)};
\node[rotate=-9, text=cfMuted!35, font=\ttfamily\footnotesize] at (19.0,24.9) {$\mathsf{Tuple}\to\mathsf{Nest}$};
\node[rotate=6, text=cfMuted!32, font=\ttfamily\footnotesize] at (2.2,20.4) {$A\circ B$};
\node[rotate=-6, text=cfMuted!32, font=\ttfamily\footnotesize] at (19.3,18.7) {logical\_divide};
\node[rotate=-7, text=cfMuted!30, font=\ttfamily\footnotesize] at (2.6,6.5) {((2,2),(2,2))};
\node[rotate=7, text=cfMuted!30, font=\ttfamily\footnotesize] at (19.0,6.9) {size : cosize};

% 顶部小标签
\node[text=cfMuted, font=\ttfamily\bfseries\small] at (10.8,25.7)
  {CATEGORY\;THEORY $\cdot$ CUTE\;LAYOUTS $\cdot$ ARXIV:2601.05972};

% 标题块
\node[text=cfInk, font=\fontsize{46}{50}\selectfont\bfseries] at (10.8,22.9) {CuTe 布局};
\node[text=cfInk, font=\fontsize{30}{36}\selectfont\bfseries] at (10.8,20.8) {范畴论基础};
\draw[cfViolet, line width=2.1pt, sketch] (5.5,20.0) -- (16.1,20.0);
\draw[cfDeep!72, line width=1.1pt, sketch] (6.0,19.6) -- (15.6,19.6);
\node[text=cfMuted, font=\fontsize{13.5}{17}\selectfont] at (10.8,18.7)
  {$\mathsf{Tuple}\;\cdot\;\mathsf{Nest}\;\cdot\;$ 可处理布局 $\cdot$ 复合 $\cdot$ 逻辑乘积 $\cdot$ 逻辑切分};

% 右上角落款式标签
\node[fill=cfDeep, text=white, rounded corners=3pt, font=\bfseries\fontsize{13}{15}\selectfont,
  inner sep=4.5pt] at (18.35,22.5) {中文译注版};

% 卡片 A（左，旋转 -5）：范畴 Tuple 的三角交换图
\begin{scope}[shift={(4.75,14.55)}, rotate=-5]
  \node[card, minimum width=6.2cm, minimum height=5.6cm] (cardA) {};
  \node[text=cfInk, font=\bfseries\small] at (0,2.25) {范畴 $\mathsf{Tuple}$};
  \node[text=cfMuted, font=\scriptsize] at (0,1.82) {对象 / 态射 / 恒等 / 复合};
  \node[obj] (a) at (-1.75,0.62) {$\langle m\rangle$};
  \node[obj] (b) at (1.65,0.62) {$\langle n\rangle$};
  \node[obj] (c) at (1.65,-1.05) {$\langle p\rangle$};
  \draw[-{Stealth[length=2mm]}, cfViolet, line width=0.9pt, sketch]
    (a) -- node[above, font=\tiny, text=cfInk] {$f$} (b);
  \draw[-{Stealth[length=2mm]}, cfViolet, line width=0.9pt, sketch]
    (b) -- node[right, font=\tiny, text=cfInk] {$g$} (c);
  \draw[-{Stealth[length=2mm]}, cfViolet, line width=0.9pt, sketch]
    (a) -- node[below left, font=\tiny, text=cfInk] {$g\circ f$} (c);
  \node[text=cfMuted!85, font=\ttfamily\scriptsize] at (0,-2.25) {$\mathsf{Tuple}\to\mathsf{Layout}$};
\end{scope}

% 卡片 B（右，旋转 4）：布局代数的三种运算
\begin{scope}[shift={(16.55,14.35)}, rotate=4]
  \node[card, minimum width=6.2cm, minimum height=5.6cm] (cardB) {};
  \node[text=cfInk, font=\bfseries\small] at (0,2.25) {布局代数};
  \node[text=cfMuted, font=\scriptsize] at (0,1.82) {与范畴运算相容};
  \node[text=cfDeep, font=\ttfamily\scriptsize, align=left] at (0,0.35)
    {comp(A, B)\\[2pt] product(A, B)\\[2pt] divide(A, B)};
  \draw[coral, line width=1.6pt, sketch] (-2.2,-1.35) -- (-1.3,-1.35);
  \node[text=cfMuted!80, font=\scriptsize] at (0,-2.05) {coalesce / complement / compact};
\end{scope}

% 卡片 C（中下，旋转 -2）：可处理布局与 colorx 八色网格
\begin{scope}[shift={(10.8,8.35)}, rotate=-2]
  \node[card, minimum width=7.6cm, minimum height=6.0cm] (cardC) {};
  \node[text=cfInk, font=\bfseries\small] at (0,2.42) {可处理布局};
  \node[text=cfMuted, font=\scriptsize] at (0,1.99) {TRACTABLE LAYOUTS};
  \foreach \x/\clr in {0/colorx1,1/colorx3,2/colorx5,3/colorx7}
    \foreach \y/\tint in {0/20,1/34,2/48,3/62}
      \fill[\clr!\tint] (-3.35+\x*0.68,-1.55+\y*0.68) rectangle +(0.60,0.60);
  \node[text=cfInk, font=\ttfamily\scriptsize] at (1.15,1.05) {L = (4,4):(1,4)};
  \node[text=cfMuted, font=\ttfamily\scriptsize] at (1.15,0.5) {size 16 $\cdot$ cosize 16};
  \draw[mint!80!black, line width=1.3pt, sketch] (-0.3,-0.05) -- (2.9,-0.05);
  \node[text=cfMuted, font=\ttfamily\scriptsize, align=left] at (1.15,-0.85)
    {compact\\[2pt] coalesce\\[2pt] complement};
  \node[text=cfMuted!80, font=\scriptsize] at (0,-2.35) {Colfax Research $\cdot$ arXiv:2601.05972};
\end{scope}

% 卡片间手绘箭头（起止点均在卡片包围盒外）
\draw[-{Stealth[length=3.2mm]}, cfDeep!70, line width=1.3pt, sketch]
  (8.05,13.55) to[bend left=14] (9.55,11.85);
\draw[-{Stealth[length=3.2mm]}, cfDeep!70, line width=1.3pt, sketch]
  (13.3,13.35) to[bend right=14] (11.95,11.85);

% 底部：colorx 八段色条 + 落款
\foreach \k/\clr in {0/colorx1,1/colorx2,2/colorx3,3/colorx4,4/colorx5,5/colorx6,6/colorx7,7/colorx8}
  \fill[\clr!85] (3.35+\k*1.86,4.9) rectangle +(1.86,0.08);
\node[text=cfInk!85, font=\fontsize{15}{18}\selectfont\bfseries] at (10.8,4.0)
  {忠实译注 $\cdot$ 术语统一 $\cdot$ 仅供学习交流};
\node[text=cfMuted, font=\small] at (10.8,3.05)
  {原文：Jack Carlisle $\cdot$ Jay Shah $\cdot$ Reuben Stern $\cdot$ Paul VanKoughnett};
\node[text=cfMuted, font=\small] at (10.8,2.45)
  {译注：LeetCUDA Project $\cdot$ \builddate};

\end{tikzpicture}
\end{document}
